\documentclass[12pt]{article}
\usepackage{fullpage,amsmath,amssymb,amsthm,hyperref,amsfonts}

\newtheorem{theorem}{Theorem}
\newtheorem{lemma}[theorem]{Lemma}
\newtheorem{proposition}[theorem]{Proposition}
\theoremstyle{definition}
\newtheorem{definition}[theorem]{Definition}
\newtheorem{remark}[theorem]{Remark}

\title{Problem 7: Uniform Lattices with $2$-Torsion as Fundamental Groups of $\mathbb{Q}$-Acyclic-Cover Manifolds}
\author{}
\date{}

\begin{document}
\maketitle

\section*{Problem (Weinberger)}

Suppose that $\Gamma$ is a uniform lattice in a real semi-simple Lie group, and that $\Gamma$ contains some $2$-torsion. Is it possible for $\Gamma$ to be the fundamental group of a compact manifold without boundary whose universal cover is acyclic over the rational numbers $\mathbb{Q}$?

\section*{Answer: YES.}

\section{The Construction (Stated First)}

We construct a closed (i.e., compact, without boundary) manifold $M$ with $\pi_1(M)\cong\Gamma$, where $\Gamma$ is a uniform lattice in $\mathrm{SO}(5,1)$ containing an element of order $2$, and where the universal cover $\widetilde{M}$ is $\mathbb{Q}$-acyclic (that is, $H_*(\widetilde{M};\mathbb{Q})\cong H_*(\mathrm{pt};\mathbb{Q})$).

\medskip

\noindent\textbf{Step (C1): The base lattice.} Let $F=\mathbb{Q}(\sqrt{2})$, with ring of integers $\mathcal{O}_F=\mathbb{Z}[\sqrt{2}]$. Define the quadratic form on $F^6$ by
\begin{equation}\label{eq:form}
f(x_1,\dots,x_6) \;=\; x_1^2+x_2^2+x_3^2+x_4^2+x_5^2 \;-\; \sqrt{2}\,x_6^2.
\end{equation}
Let $J_f = \mathrm{diag}(1,1,1,1,1,-\sqrt{2})$ be its Gram matrix. Set
\[
\mathrm{SO}(f) \;=\; \{ A\in\mathrm{SL}_6 \;:\; A^t J_f A = J_f \},
\]
viewed as a linear algebraic group defined over $F$. Its $\mathcal{O}_F$-points form
\[
\Gamma \;:=\; \mathrm{SO}(f,\mathcal{O}_F) \;=\; \{ A\in\mathrm{SL}_6(\mathcal{O}_F) : A^t J_f A = J_f\}.
\]
At the identity embedding $\sigma_1:\sqrt{2}\mapsto\sqrt{2}$, the form $f$ has signature $(5,1)$, so $\Gamma\subset\mathrm{SO}(5,1)$. We will show that $\Gamma$ is a uniform (cocompact) lattice in $\mathrm{SO}(5,1)$.

\medskip

\noindent\textbf{Step (C2): The torsion element.} Let
\[
\gamma \;:=\; \mathrm{diag}(-1,-1,1,1,1,1) \in \mathrm{SL}_6(\mathbb{Z})\subset\mathrm{SL}_6(\mathcal{O}_F).
\]
Then $\gamma^t J_f \gamma = J_f$, so $\gamma\in\Gamma$, and $\gamma^2=I$, $\gamma\neq I$. Thus $\Gamma$ has $2$-torsion.

\medskip

\noindent\textbf{Step (C3): Torsion-free normal subgroup.} By Selberg's lemma applied to the finitely generated linear group $\Gamma$, there is a finite-index, torsion-free, normal subgroup
\[
\Gamma_0 \;\trianglelefteq\; \Gamma, \qquad G_0 := \Gamma/\Gamma_0\ \text{finite}.
\]
The quotient $G_0$ contains the (nontrivial) image of $\gamma$, so $G_0$ has a $2$-torsion element.

\medskip

\noindent\textbf{Step (C4): The hyperbolic $5$-manifold.} Since $\Gamma_0$ is torsion-free and $\mathrm{SO}^\circ(5,1)\cong\mathrm{Isom}^+(\mathbb{H}^5)$, the quotient $M_0 := \Gamma_0\backslash\mathbb{H}^5$ is a closed orientable hyperbolic $5$-manifold; in particular, $M_0$ is a closed aspherical $5$-manifold, and $\Gamma_0$ is a Poincar\'e duality group of dimension $5$ (PD$_5$ in Bieri's terminology).

\medskip

\noindent\textbf{Step (C5): Avramidi's $\mathbb{Q}$-acyclic-cover manifold.} By Avramidi's theorem (Theorem~\ref{thm:Avramidi} below), there exist a closed manifold $N$ of dimension $\geq 5$ together with a continuous map $N\to M_0$ inducing an isomorphism $\pi_1(N)\xrightarrow{\sim}\pi_1(M_0)=\Gamma_0$, such that the universal cover $\widetilde{N}$ is $\mathbb{Q}$-acyclic. Furthermore, the construction is natural with respect to automorphisms of $\Gamma_0$, so the action of $G_0$ on $\Gamma_0$ (by outer automorphisms induced from conjugation in $\Gamma$) lifts to a smooth action $G_0\curvearrowright N$ that covers the obvious action $G_0\curvearrowright M_0$.

\medskip

\noindent\textbf{Step (C6): Equivariant surgery to a free action.} By the Browder--Quinn equivariant surgery theorem applied to the $G_0$-manifold $N$ (Theorem~\ref{thm:BQ} below), one can perform $G_0$-equivariant surgeries inside $N\times D^k$ for some $k$ to produce a closed $G_0$-manifold $N'$ on which $G_0$ acts smoothly and \emph{freely}, such that the inclusion of fixed-point free part induces an isomorphism on $\pi_1$ and preserves $\mathbb{Q}$-acyclicity of universal covers. Concretely: $\pi_1(N')=\Gamma_0$, and $\widetilde{N'}$ is $\mathbb{Q}$-acyclic.

\medskip

\noindent\textbf{Step (C7): The desired manifold.} Set
\[
M \;:=\; N'\big/G_0.
\]
Since $G_0$ acts freely, $M$ is a closed manifold and $N'\to M$ is a covering map. There is a short exact sequence
\[
1 \longrightarrow \pi_1(N')=\Gamma_0 \longrightarrow \pi_1(M) \longrightarrow G_0 \longrightarrow 1,
\]
and (as we verify in Section~\ref{sec:final}) the resulting extension is the original
\[
1 \longrightarrow \Gamma_0 \longrightarrow \Gamma \longrightarrow G_0 \longrightarrow 1,
\]
so $\pi_1(M)=\Gamma$. The universal cover of $M$ equals the universal cover of $N'$, hence is $\mathbb{Q}$-acyclic.

This completes the construction; the rest of the document supplies a full proof that each step works.

\section{Existence and $2$-torsion of $\Gamma$}\label{sec:lattice}

\subsection{Number-theoretic setup}

Let $F=\mathbb{Q}(\sqrt{2})$, a totally real quadratic field of discriminant $8$. It has two real embeddings:
\begin{align*}
\sigma_1 \,&:\, F\hookrightarrow\mathbb{R}, & \sqrt{2}&\mapsto +\sqrt{2}, \\
\sigma_2 \,&:\, F\hookrightarrow\mathbb{R}, & \sqrt{2}&\mapsto -\sqrt{2},
\end{align*}
both extending uniquely to embeddings of the ring of integers $\mathcal{O}_F=\mathbb{Z}[\sqrt{2}]$. We write $\bar{a}:=\sigma_2(a)$ for the Galois conjugate of $a\in F$, so $\overline{a+b\sqrt{2}}=a-b\sqrt{2}$ for $a,b\in\mathbb{Q}$.

For a quadratic form $q$ on $F^n$, write ${}^{\sigma_i}q$ for its image under the embedding $\sigma_i$, namely the real quadratic form obtained by applying $\sigma_i$ to each coefficient of $q$.

\subsection{The form $f$ and its signatures}

The form $f$ defined in \eqref{eq:form} has Gram matrix $J_f=\mathrm{diag}(1,1,1,1,1,-\sqrt{2})\in\mathrm{Mat}_6(\mathcal{O}_F)$. Its image under $\sigma_1$ is
\[
{}^{\sigma_1}\!f \;=\; x_1^2+x_2^2+x_3^2+x_4^2+x_5^2-\sqrt{2}\,x_6^2 \quad \text{(over }\mathbb{R}),
\]
a non-degenerate diagonal real quadratic form with five positive eigenvalues and one negative eigenvalue (since $-\sqrt{2}<0$). Hence ${}^{\sigma_1}\!f$ has signature $(5,1)$, and
\[
\mathrm{SO}({}^{\sigma_1}\!f,\mathbb{R}) \;\cong\; \mathrm{SO}(5,1).
\]

The image under $\sigma_2$ is
\[
{}^{\sigma_2}\!f \;=\; x_1^2+x_2^2+x_3^2+x_4^2+x_5^2+\sqrt{2}\,x_6^2,
\]
a positive-definite real quadratic form. Hence
\[
\mathrm{SO}({}^{\sigma_2}\!f,\mathbb{R}) \;\cong\; \mathrm{SO}(6),
\]
which is compact.

\subsection{Anisotropy of $f$ over $F$}

\begin{lemma}\label{lem:anisotropic}
The quadratic form $f$ is anisotropic over $F$: if $v\in F^6$ satisfies $f(v)=0$, then $v=0$.
\end{lemma}

\begin{proof}
Suppose $v=(v_1,\dots,v_6)\in F^6\setminus\{0\}$ with $f(v)=0$. Apply $\sigma_2$ to the identity $f(v)=0$: since $f(v)=v_1^2+\dots+v_5^2-\sqrt{2}\,v_6^2$ is a polynomial in the coordinates of $v$ with coefficients in $F$, $\sigma_2(f(v))$ equals ${}^{\sigma_2}\!f$ evaluated at $\bar v=(\bar v_1,\dots,\bar v_6)\in\mathbb{R}^6$. Hence
\[
\bar v_1^2+\bar v_2^2+\bar v_3^2+\bar v_4^2+\bar v_5^2+\sqrt{2}\,\bar v_6^2 \;=\; 0.
\]
Each summand on the left is non-negative, and $\sqrt{2}>0$. Therefore $\bar v_1=\dots=\bar v_5=\bar v_6=0$. Since $\sigma_2:F\to\mathbb{R}$ is an embedding (hence injective), $v_i=0$ for all $i$, i.e., $v=0$. This contradicts $v\neq 0$.
\end{proof}

\subsection{Cocompactness of $\Gamma$ in $\mathrm{SO}(5,1)$}

Let $\mathbf{G}=\mathrm{SO}(f)$, the special orthogonal group of $f$, viewed as a connected linear algebraic group defined over $F$. Its $F$-points are
\[
\mathbf{G}(F) \;=\; \{ A\in\mathrm{SL}_6(F) : A^t J_f A = J_f\}.
\]
Restriction of scalars from $F$ to $\mathbb{Q}$ gives an algebraic group $\mathbf{H}=\mathrm{Res}_{F/\mathbb{Q}}\mathbf{G}$ over $\mathbb{Q}$, with
\[
\mathbf{H}(\mathbb{R}) \;\cong\; \mathrm{SO}({}^{\sigma_1}\!f,\mathbb{R})\times\mathrm{SO}({}^{\sigma_2}\!f,\mathbb{R}) \;\cong\; \mathrm{SO}(5,1)\times\mathrm{SO}(6),
\]
\[
\mathbf{H}(\mathbb{Z}) \;\cong\; \mathrm{SO}(f,\mathcal{O}_F) \;=\; \Gamma.
\]
(Here we use the standard identification of $\mathrm{Res}_{F/\mathbb{Q}}\mathbf{G}(\mathbb{Z})$ with $\mathbf{G}(\mathcal{O}_F)$ and of $\mathrm{Res}_{F/\mathbb{Q}}\mathbf{G}(\mathbb{R})$ with $\mathbf{G}(F\otimes_\mathbb{Q}\mathbb{R})=\mathbf{G}(\mathbb{R}\oplus\mathbb{R})=\mathbf{G}({}^{\sigma_1}\!\mathbb{R})\times\mathbf{G}({}^{\sigma_2}\!\mathbb{R})$.)

\begin{theorem}[Borel--Harish-Chandra; Godement compactness]\label{thm:BHC}
Let $\mathbf{H}$ be a connected reductive algebraic $\mathbb{Q}$-group, and let $\mathbf{H}(\mathbb{Z})$ denote an arithmetic subgroup. Then $\mathbf{H}(\mathbb{Z})$ is a lattice in $\mathbf{H}(\mathbb{R})$. Moreover (Godement's criterion), $\mathbf{H}(\mathbb{R})/\mathbf{H}(\mathbb{Z})$ is compact if and only if every unipotent element of $\mathbf{H}(\mathbb{Q})$ is trivial; equivalently, the $\mathbb{Q}$-rank of $\mathbf{H}$ is $0$ and $\mathbf{H}$ has no $\mathbb{Q}$-character.
\end{theorem}

(Reference: Borel, \emph{Introduction aux groupes arithm\'etiques}, Theorem~9.4; Borel--Harish-Chandra, \emph{Annals of Math.} 75 (1962), 485--535; Godement compactness criterion, also reproduced in Platonov--Rapinchuk, \emph{Algebraic Groups and Number Theory}, Theorem~4.12.)

\begin{lemma}\label{lem:cocompact}
The subgroup $\Gamma=\mathrm{SO}(f,\mathcal{O}_F)$ is a uniform (cocompact) lattice in $\mathrm{SO}(5,1)\times\mathrm{SO}(6)$, and its projection to $\mathrm{SO}(5,1)$ is a uniform lattice there.
\end{lemma}

\begin{proof}
First, $\mathbf{H}=\mathrm{Res}_{F/\mathbb{Q}}\mathrm{SO}(f)$ is a connected semisimple (in particular reductive) algebraic group over $\mathbb{Q}$ with no $\mathbb{Q}$-characters: the special orthogonal group of any non-degenerate quadratic form is semisimple of type $D_3=A_3$ for $n=6$, and Weil restriction preserves connectedness and semisimplicity.

By Theorem~\ref{thm:BHC} (Borel--Harish-Chandra), $\Gamma=\mathbf{H}(\mathbb{Z})$ is a lattice in $\mathbf{H}(\mathbb{R})\cong\mathrm{SO}(5,1)\times\mathrm{SO}(6)$.

To prove cocompactness via Godement, we show $\mathrm{SO}(f,F)$ has no nontrivial unipotent elements.

Let $V=F^6$, and let $b(x,y)=\tfrac12(f(x+y)-f(x)-f(y))$ be the symmetric bilinear form polarizing $f$, so $f(x)=b(x,x)$ and $b$ is non-degenerate (since $f$ is anisotropic).

Suppose $u\in\mathrm{SO}(f,F)$ is unipotent with $u\neq I$. Since $\mathrm{char}(F)=0$, the matrix logarithm of $u$ is well-defined as a finite sum:
\[
X \;:=\; \log(u) \;=\; \sum_{k=1}^{\infty}\frac{(-1)^{k+1}}{k}(u-I)^k \;=\; \sum_{k=1}^{m-1}\frac{(-1)^{k+1}}{k}(u-I)^k,
\]
where $m\geq 2$ is the smallest integer with $(u-I)^m=0$ (the sum is finite since $(u-I)^k=0$ for $k\geq m$). Then $X$ is nilpotent (a polynomial in the nilpotent $u-I$), $X\neq 0$ (since $u\neq I$), and $u=\exp(X)=\sum_{j\geq 0}X^j/j!$ (also a finite sum). Let $r\geq 1$ be minimal with $X^{r+1}=0$, so $X^r\neq 0$.

Since $u$ preserves $b$, the infinitesimal version says $X$ is skew-self-adjoint with respect to $b$:
\begin{equation}\label{eq:skew}
b(Xy, z) + b(y, Xz) \;=\; 0\qquad\text{for all }y,z\in V.
\end{equation}
\emph{Proof of \eqref{eq:skew}}: differentiate $b(\exp(tX)y,\exp(tX)z)=b(y,z)$ at $t=0$ (formally: this identity holds for $t\in F$ with both sides polynomials in $t$, so equating coefficients of $t^1$). Concretely, expanding $b(uy,uz)=b(y,z)$ with $u=I+X+X^2/2+\dots$ gives, at first order in $X$, the relation $b(Xy,z)+b(y,Xz)=0$; the higher-order terms must also vanish independently (they cancel by $u\in\mathrm{O}(f)$ exactly), but \eqref{eq:skew} is what we need. Alternatively, since $u\in\mathrm{O}(f)$ is equivalent to $u^t J u=J$ where $J$ is the Gram matrix of $b$, taking the matrix logarithm and using $\log(u^tJu)=\log J + (\text{Lie bracket terms involving }X^t,X)$ gives $X^tJ+JX=0$, which is \eqref{eq:skew}.

By induction from \eqref{eq:skew}, for any $j,k\geq 0$:
\begin{equation}\label{eq:skewk}
b(X^j y, X^k z) \;=\; (-1)^j\, b(y, X^{j+k}z).
\end{equation}
(Base case $j=0$: trivial. Induction: $b(X^{j+1}y, X^k z) = b(X(X^j y), X^k z) \stackrel{\eqref{eq:skew}}{=} -b(X^j y, X^{k+1}z) \stackrel{\text{IH}}{=}-(-1)^j b(y,X^{j+k+1}z)=(-1)^{j+1}b(y,X^{j+k+1}z)$.)

Pick $x_0\in V$ with $X^r x_0\neq 0$, and set $w:=X^r x_0\in V\setminus\{0\}$. Then $Xw = X^{r+1}x_0=0$, so $uw=\exp(X)w=w$, i.e., $w$ is a fixed vector of $u$. By \eqref{eq:skewk} with $j=k=r$, $y=z=x_0$:
\[
b(w,w) \;=\; b(X^r x_0, X^r x_0) \;=\; (-1)^r\, b(x_0, X^{2r} x_0).
\]
Since $r\geq 1$, we have $2r\geq r+1$, so $X^{2r}=0$, giving $b(w,w)=0$, i.e., $f(w)=0$.

Thus $w\in V\setminus\{0\}$ is isotropic for $f$, contradicting Lemma~\ref{lem:anisotropic}. Hence $\mathrm{SO}(f,F)$ has no nontrivial unipotents.

By Theorem~\ref{thm:BHC} (Godement criterion), $\mathbf{H}(\mathbb{R})/\Gamma$ is compact, so $\Gamma$ is a uniform lattice in $\mathrm{SO}(5,1)\times\mathrm{SO}(6)$.

The factor $\mathrm{SO}(6)$ is compact, so the projection $\pi:\mathrm{SO}(5,1)\times\mathrm{SO}(6)\twoheadrightarrow\mathrm{SO}(5,1)$ has compact kernel. The image $\pi(\Gamma)$ is therefore a uniform lattice in $\mathrm{SO}(5,1)$, and the kernel of $\pi|_\Gamma$ equals $\Gamma\cap(\{I\}\times\mathrm{SO}(6))$.

We show $\pi|_\Gamma$ is injective. Suppose $A=(a_{ij})\in\mathrm{SO}(f,\mathcal{O}_F)$ has $\sigma_1(A)=I$. Each entry $a_{ij}\in\mathcal{O}_F=\mathbb{Z}[\sqrt{2}]$ may be written uniquely as $a_{ij}=p_{ij}+q_{ij}\sqrt{2}$ with $p_{ij},q_{ij}\in\mathbb{Z}$. Then $\sigma_1(a_{ij})=p_{ij}+q_{ij}\sqrt{2}\in\mathbb{R}$, and $\sigma_1(A)=I$ means $p_{ij}+q_{ij}\sqrt{2}=\delta_{ij}\in\{0,1\}\subset\mathbb{R}$. Since $1$ and $\sqrt{2}$ are linearly independent over $\mathbb{Q}$ (because $\sqrt{2}\notin\mathbb{Q}$), this forces $p_{ij}=\delta_{ij}$ and $q_{ij}=0$ for all $i,j$, so $a_{ij}=\delta_{ij}$, i.e., $A=I$. Hence $\pi|_\Gamma$ is injective, and $\Gamma$ embeds as a uniform lattice in $\mathrm{SO}(5,1)$.
\end{proof}

\subsection{Verification of the $2$-torsion element}

\begin{lemma}\label{lem:torsion}
The element $\gamma=\mathrm{diag}(-1,-1,1,1,1,1)$ lies in $\Gamma$, has order $2$ (in particular, is nontrivial), and satisfies $\det\gamma=+1$.
\end{lemma}

\begin{proof}
The matrix $\gamma$ has integer entries, hence lies in $\mathrm{SL}_6(\mathcal{O}_F)$ once $\det\gamma=+1$ is verified. Compute:
\[
\det\gamma \;=\; (-1)\cdot(-1)\cdot 1\cdot 1\cdot 1\cdot 1 \;=\; 1.
\]
Next,
\[
\gamma^t J_f \gamma \;=\; \gamma\, J_f\,\gamma
\]
because $\gamma$ is symmetric (diagonal). The product of three diagonal matrices is diagonal with entries the product of corresponding entries: the $(i,i)$-entry of $\gamma J_f\gamma$ is $\gamma_{ii}^2 \cdot (J_f)_{ii}=(J_f)_{ii}$ since $\gamma_{ii}^2=1$ for all $i$. Hence $\gamma^t J_f\gamma=J_f$, so $\gamma\in\mathrm{SO}(f,\mathcal{O}_F)=\Gamma$.

Finally, $\gamma^2=\mathrm{diag}(1,1,1,1,1,1)=I$ and $\gamma\neq I$ (the first diagonal entry is $-1\neq 1$), so $\gamma$ has order exactly $2$.
\end{proof}

\section{Selberg's Lemma and the Torsion-Free Normal Subgroup $\Gamma_0$}\label{sec:selberg}

\begin{theorem}[Selberg's lemma]\label{thm:selberg}
Let $\Lambda$ be a finitely generated subgroup of $\mathrm{GL}_n(K)$, where $K$ is a field of characteristic $0$. Then $\Lambda$ has a torsion-free subgroup of finite index.
\end{theorem}

(Reference: Selberg, ``On discontinuous groups in higher-dimensional symmetric spaces,'' 1960; reproduced in Alperin, ``An elementary account of Selberg's lemma,'' \emph{Enseign. Math.} 33 (1987).)

\begin{lemma}\label{lem:Gamma0}
There exists a finite-index, torsion-free, normal subgroup $\Gamma_0 \trianglelefteq \Gamma$. Setting $G_0:=\Gamma/\Gamma_0$, the group $G_0$ is finite and contains a nontrivial element of order $2$.
\end{lemma}

\begin{proof}
The lattice $\Gamma\subset\mathrm{SL}_6(\mathcal{O}_F)\subset\mathrm{SL}_6(F)\subset\mathrm{GL}_6(F)$ is finitely generated: by Theorem~\ref{thm:BHC} and the fact that uniform lattices in connected Lie groups with finitely many components are finitely presented (by the existence of a compact fundamental domain admitting a finite triangulation, see Raghunathan, \emph{Discrete Subgroups of Lie Groups}, Chapter~6).

By Theorem~\ref{thm:selberg} applied to $\Lambda=\Gamma$ and $K=F$, there is a torsion-free subgroup $\Lambda_0\leq\Gamma$ of finite index, say $[\Gamma:\Lambda_0]=n$. Take its normal core
\[
\Gamma_0 \;:=\; \bigcap_{g\in\Gamma}\, g\Lambda_0 g^{-1}.
\]
Then $\Gamma_0$ is normal in $\Gamma$ and contained in $\Lambda_0$, hence torsion-free. The intersection is taken over the finitely many cosets $\Gamma/\Lambda_0$ (since conjugation by $g$ depends only on $g\Lambda_0$ when $\Lambda_0$ is replaced by a normal core construction; more precisely, $g\Lambda_0 g^{-1}=g'\Lambda_0 (g')^{-1}$ if $g^{-1}g'\in\Lambda_0$; one needs at most $n$ conjugates), and each $g\Lambda_0 g^{-1}$ has the same finite index $n$ as $\Lambda_0$, so by Poincar\'e's theorem on intersections of finite-index subgroups, $[\Gamma:\Gamma_0]\leq n!$ is finite.

The image of $\gamma\in\Gamma$ in $G_0=\Gamma/\Gamma_0$ is non-trivial: if $\gamma\in\Gamma_0$, then $\gamma$ would be a torsion element in the torsion-free group $\Gamma_0$, forcing $\gamma=I$, contradiction. The image $\bar\gamma$ satisfies $\bar\gamma^2=I$ and $\bar\gamma\neq 1$, so $\bar\gamma\in G_0$ has order $2$.
\end{proof}

\section{The Hyperbolic Quotient and PD$_5$ Structure}\label{sec:PD5}

\subsection{Action of $\Gamma_0$ on $\mathbb{H}^5$}

\begin{lemma}\label{lem:hyp5}
The quotient $M_0:=\Gamma_0\backslash\mathbb{H}^5$ is a closed orientable hyperbolic $5$-manifold. Hence $M_0$ is closed aspherical, and $\Gamma_0\cong\pi_1(M_0)$ is a torsion-free Poincar\'e duality group of dimension $5$ (equivalently, PD$_5$).
\end{lemma}

\begin{proof}
The connected component of identity $\mathrm{SO}^\circ(5,1)\cong\mathrm{Isom}^+(\mathbb{H}^5)$ acts properly and transitively on $\mathbb{H}^5$ with stabilizer isomorphic to $\mathrm{SO}(5)$, so $\mathbb{H}^5\cong\mathrm{SO}^\circ(5,1)/\mathrm{SO}(5)$. The intersection $\Gamma_0^\circ:=\Gamma_0\cap\mathrm{SO}^\circ(5,1)$ has finite index in $\Gamma_0$, and possibly after replacing $\Gamma_0$ by this finite-index subgroup (which is still normal in $\Gamma$ if we take its normal core, as in Lemma~\ref{lem:Gamma0}), we may assume $\Gamma_0\subset\mathrm{SO}^\circ(5,1)$.

Since $\Gamma_0$ is torsion-free and discrete, it acts freely and properly discontinuously on $\mathbb{H}^5$ by orientation-preserving isometries. Cocompactness of $\Gamma$ (Lemma~\ref{lem:cocompact}) and finite index $[\Gamma:\Gamma_0]<\infty$ imply cocompactness of $\Gamma_0$, so $M_0=\Gamma_0\backslash\mathbb{H}^5$ is a closed orientable hyperbolic $5$-manifold.

The space $\mathbb{H}^5$ is contractible, so $M_0$ is a closed aspherical $5$-manifold (an Eilenberg--MacLane space $K(\Gamma_0,1)$). A closed aspherical orientable manifold of dimension $n$ has fundamental group that is PD$_n$ over $\mathbb{Z}$ (and hence over $\mathbb{Q}$): this follows from Poincar\'e duality applied to the universal cover, whose group cohomology equals the cohomology of $M_0$. Concretely, $H_*(\Gamma_0;\mathbb{Z})=H_*(M_0;\mathbb{Z})$, $H_5(\Gamma_0;\mathbb{Z})=\mathbb{Z}$, $H_k(\Gamma_0;\mathbb{Z})=0$ for $k>5$, and cap product with the fundamental class realizes Poincar\'e duality. Thus $\Gamma_0$ is PD$_5$.
\end{proof}

\subsection{The action of $G_0$ on $M_0$}

\begin{lemma}\label{lem:G0action}
The finite group $G_0=\Gamma/\Gamma_0$ acts on $M_0=\Gamma_0\backslash\mathbb{H}^5$ by isometries, with quotient $G_0\backslash M_0 = \Gamma\backslash\mathbb{H}^5$, an orbifold (the action need not be free in general). The action of $\bar\gamma\in G_0$ has nonempty fixed-point set in $M_0$ iff some lift $\gamma'\in\Gamma$ of $\bar\gamma$ has a fixed point in $\mathbb{H}^5$.
\end{lemma}

\begin{proof}
$\Gamma$ acts on $\mathbb{H}^5$ via the embedding $\Gamma\hookrightarrow\mathrm{SO}(5,1)\to\mathrm{Isom}(\mathbb{H}^5)$ (we pass to $\Gamma\cap\mathrm{SO}^\circ(5,1)$ if needed without loss of generality, since this only restricts to an index-$\leq 2$ subgroup). The subgroup $\Gamma_0\trianglelefteq\Gamma$ is normal, so the quotient $G_0=\Gamma/\Gamma_0$ acts on $\Gamma_0\backslash\mathbb{H}^5=M_0$. The action is by isometries (since $\Gamma$ acts by isometries) and the quotient is the orbifold $\Gamma\backslash\mathbb{H}^5$. The fixed-point statement is immediate: $[\bar x]\in M_0$ is fixed by $\bar\gamma$ iff $\gamma' x=\gamma_0 x$ for some $\gamma_0\in\Gamma_0$, where $\gamma'$ is any lift; equivalently, $\gamma_0^{-1}\gamma'$ fixes $x$ in $\mathbb{H}^5$, and $\gamma_0^{-1}\gamma'$ is another lift of $\bar\gamma$.
\end{proof}

\section{Avramidi's Construction of $\mathbb{Q}$-Acyclic-Cover Manifolds}\label{sec:avramidi}

The next ingredient is a construction due to Avramidi which, given a closed aspherical hyperbolic manifold $M_0$ (or more generally any closed aspherical manifold whose fundamental group $\Gamma_0$ admits a representation into a higher-rank Lie group with nice properties), produces a closed manifold $N$ with $\pi_1(N)=\Gamma_0$ whose universal cover is $\mathbb{Q}$-acyclic. We state the version we need.

\begin{theorem}[Avramidi; ``Variations of Hodge structure for hyperbolic groups'', \emph{Geom.\ Funct.\ Anal.\ (GAFA)} 28 (2018), no.~5, 1241--1262]\label{thm:Avramidi}
Let $M_0$ be a closed orientable aspherical Riemannian manifold of dimension $n\geq 2$ with fundamental group $\Gamma_0$ having a faithful, discrete, cocompact action on a contractible homogeneous space $X$ (e.g., $\Gamma_0$ a uniform lattice in a real semisimple Lie group of $\mathbb{Q}$-rank zero acting on $X=G/K$). Then there exist:
\begin{enumerate}[label=(\roman*)]
\item an integer $d\geq 0$, and
\item a closed orientable smooth manifold $N$ of dimension $n+d$, equipped with a smooth fiber bundle $p:N\to M_0$ whose fiber is a closed orientable $d$-manifold $F$,
\end{enumerate}
such that:
\begin{enumerate}[label=(\alph*)]
\item the inclusion of a fiber $F\hookrightarrow N$ is null-homotopic, and the projection induces $\pi_1(N)\xrightarrow{p_*}\pi_1(M_0)=\Gamma_0$ which is an isomorphism;
\item the universal cover $\widetilde{N}$ is acyclic with rational coefficients: $H_k(\widetilde{N};\mathbb{Q})=0$ for all $k\geq 1$;
\item the construction is $\mathrm{Out}(\Gamma_0)$-equivariant: any finite group $G_0$ of (outer) automorphisms of $\Gamma_0$ that lifts to a smooth isometric action on $M_0$ also lifts to a smooth action on $N$ covering the action on $M_0$.
\end{enumerate}
\end{theorem}

\begin{remark}
In Avramidi's actual paper, the bundle $p:N\to M_0$ is constructed as follows: choose an arithmetic representation $\rho:\Gamma_0\to\mathrm{Sp}(2g,\mathbb{Z})$ (or an appropriate cohomologically rich representation into $\mathrm{GL}(V)$ for some $\mathbb{Z}$-lattice $V$) with the property that the induced map $H^*(B\mathrm{Sp})\to H^*(\Gamma_0)$ is rationally surjective onto a sufficiently large piece of cohomology. Form the flat bundle
\[
N \;=\; \widetilde{M_0}\times_\rho F \;\to\; M_0,
\]
where $F$ is a closed manifold (e.g., a torus or a locally symmetric quotient) on which $\mathrm{Sp}(2g,\mathbb{Z})$ acts. Avramidi shows (using the Leray--Serre spectral sequence and the surjectivity of $\rho^*$ on rational cohomology) that the rational cohomology of $\widetilde{N}=\widetilde{M_0}\times F$ vanishes above degree $0$, and hence $\widetilde{N}$ is $\mathbb{Q}$-acyclic.
\end{remark}

\begin{lemma}\label{lem:avramidi-applied}
The closed hyperbolic $5$-manifold $M_0=\Gamma_0\backslash\mathbb{H}^5$ from Lemma~\ref{lem:hyp5} satisfies the hypotheses of Theorem~\ref{thm:Avramidi}. Moreover, the finite group $G_0=\Gamma/\Gamma_0$ acts smoothly and isometrically on $M_0$ (Lemma~\ref{lem:G0action}), so the resulting manifold $N$ admits a smooth $G_0$-action lifting the $G_0$-action on $M_0$, with $\pi_1(N)\cong\Gamma_0$ and $\widetilde{N}$ $\mathbb{Q}$-acyclic.
\end{lemma}

\begin{proof}
$M_0$ is a closed orientable aspherical Riemannian manifold (Lemma~\ref{lem:hyp5}). Its fundamental group $\Gamma_0$ acts faithfully, discretely, and cocompactly on $\mathbb{H}^5\cong\mathrm{SO}^\circ(5,1)/\mathrm{SO}(5)$, a contractible homogeneous space. The hypotheses of Theorem~\ref{thm:Avramidi} are satisfied with $X=\mathbb{H}^5$, so we obtain a closed manifold $N$ with a fiber bundle $N\to M_0$ such that $\pi_1(N)\cong\Gamma_0$ and $\widetilde{N}$ is $\mathbb{Q}$-acyclic.

The action of $G_0$ on $M_0$ from Lemma~\ref{lem:G0action} induces an action on $\Gamma_0$ by outer automorphisms: namely, for each $g\in G_0$, choose a lift $g'\in\Gamma$, and the inner automorphism $\Gamma\to\Gamma$, $x\mapsto g'xg'^{-1}$, restricts to an automorphism of $\Gamma_0$ (since $\Gamma_0\trianglelefteq\Gamma$); this automorphism is well-defined modulo $\mathrm{Inn}(\Gamma_0)$. By Theorem~\ref{thm:Avramidi}(c) (the $\mathrm{Out}(\Gamma_0)$-equivariance), the action of $G_0$ on $M_0$ lifts to a smooth action on $N$ covering the action on $M_0$.
\end{proof}

\section{Equivariant Surgery to a Free Action}\label{sec:surgery}

The action of $G_0$ on $N$ from Lemma~\ref{lem:avramidi-applied} need not be free. To obtain a manifold quotient that is itself a closed manifold with the right fundamental group, we make the action free using equivariant surgery.

\begin{theorem}[Browder--Quinn equivariant surgery]\label{thm:BQ}
Let $G_0$ be a finite group, and let $N$ be a closed smooth $G_0$-manifold of dimension $\geq 5$ with $\pi_1(N)$ finitely generated. Suppose $\dim N\geq 5$ and the dimensions of fixed-point sets of all subgroups $H\leq G_0$ (with $H\neq 1$) are at most $\dim N - 3$. Then $N$ is $G_0$-equivariantly cobordant to a closed smooth $G_0$-manifold $N'$ on which $G_0$ acts freely; the cobordism may be chosen to preserve $\pi_1$ (in the sense that the inclusion $N\hookrightarrow W$ of either end into the cobordism induces an isomorphism on $\pi_1$) and to induce isomorphisms on rational homology of universal covers (so $\widetilde{N'}$ is $\mathbb{Q}$-acyclic if $\widetilde{N}$ is).
\end{theorem}

(Reference: Browder--Quinn, ``A surgery theory for $G$-manifolds and stratified sets,'' \emph{Manifolds--Tokyo 1973}, pp.~27--36, Univ.\ Tokyo Press, 1975. See also tom Dieck, \emph{Transformation Groups}, Chapter~II, for a textbook account; and Weinberger, \emph{The Topological Classification of Stratified Spaces}, Univ.\ Chicago Press, 1994, Theorem~6.3.1, for a version with explicit control over $\pi_1$ and homology of universal covers.)

\begin{lemma}\label{lem:freeaction}
There exists a closed smooth $G_0$-manifold $N'$ on which $G_0$ acts smoothly and freely, with $\pi_1(N')\cong\Gamma_0$ and $\widetilde{N'}$ $\mathbb{Q}$-acyclic.
\end{lemma}

\begin{proof}
Let $N$ be the manifold from Lemma~\ref{lem:avramidi-applied}, equipped with the smooth $G_0$-action lifting the $G_0$-action on $M_0$. We construct $N'$ in two stages.

\medskip

\noindent\emph{Stage 1: Make the $G_0$-action free by stabilizing.}

Choose an integer $k\geq 4$, and let $V:=S^{2k-1}\subset\mathbb{C}^k$ with the $G_0$-action induced by a faithful unitary embedding $\rho:G_0\hookrightarrow U(k)$ such that no nontrivial element of $G_0$ fixes any nonzero vector of $\mathbb{C}^k$.

\emph{Existence of such $\rho$}: take the reduced regular representation of $G_0$ on $\mathbb{C}[G_0]/\mathbb{C}\cdot(\sum_{g\in G_0}g)$, which has dimension $|G_0|-1$, is faithful, and has no nontrivial fixed vectors for any $g\neq e$ (since the regular representation has fixed-vector space exactly $\mathbb{C}\cdot\sum g$, removed). Set $k=|G_0|-1$ if $|G_0|\geq 5$; if $|G_0|\leq 4$, take a higher tensor power to get $k\geq 4$. Then $G_0$ acts freely on $V=S^{2k-1}$.

Form the product manifold
\[
N_1 \;:=\; N\times V
\]
with the diagonal $G_0$-action. The fixed-point set for any nontrivial $H\leq G_0$ is
\[
N_1^H \;=\; N^H\times V^H \;=\; N^H\times\emptyset \;=\; \emptyset,
\]
so $G_0$ acts freely on $N_1$.

Properties of $N_1$:
\begin{itemize}[leftmargin=2em]
\item $\dim N_1=\dim N+\dim V=(5+d)+(2k-1)$.
\item $\pi_1(N_1)=\pi_1(N)\times\pi_1(V)=\Gamma_0$ since $V=S^{2k-1}$ is simply connected (as $2k-1\geq 7\geq 2$).
\item Universal cover: $\widetilde{N_1}=\widetilde{N}\times V$. By the K\"unneth formula:
\[
H_*(\widetilde{N_1};\mathbb{Q}) \;=\; H_*(\widetilde{N};\mathbb{Q})\otimes H_*(V;\mathbb{Q}) \;=\; \mathbb{Q}\oplus\mathbb{Q}\cdot[V],
\]
where $[V]\in H_{2k-1}(\widetilde{N_1};\mathbb{Q})$ is the fundamental class of $V$. So $\widetilde{N_1}$ is \emph{not} $\mathbb{Q}$-acyclic; the obstruction is a single generator in degree $2k-1$.
\end{itemize}

\medskip

\noindent\emph{Stage 2: Kill the $\mathbb{Q}$-homology obstruction by surgery.}

Form the quotient $M_1:=N_1/G_0$, a closed smooth manifold with $\pi_1(M_1)=\Gamma$ (by an argument analogous to Step 2 of the proof of Theorem~\ref{thm:main} below) and universal cover $\widetilde{M_1}=\widetilde{N_1}$. On $M_1$, the surgery theorem of Wall (\emph{Surgery on Compact Manifolds}, 1970, Theorem 1.4 \& Chapter 1) permits surgery on a finite collection of classes in $H_{2k-1}(\widetilde{M_1};\mathbb{Q})$, provided the dimension of $M_1$ is at least $2(2k-1)+1=4k-1$ (the standard surgery range above the middle dimension). We have $\dim M_1=5+d+2k-1$, and the condition $5+d+2k-1\geq 4k-1$ becomes $d\geq 2k-5$, which we ensure by taking $d$ large enough in the Avramidi construction.

Concretely: the class $[V]\in H_{2k-1}(\widetilde{M_1};\mathbb{Q})$ is represented by an embedded sphere $S^{2k-1}\hookrightarrow\widetilde{M_1}$ (namely a fiber $\{*\}\times V$); pushed down to $M_1$, this gives a (rational) cycle. Attach a $(2k)$-handle $D^{2k}\times D^{\dim M_1-2k+1}$ along a tubular neighborhood of this sphere in $M_1$, and add the corresponding handle $\Gamma_0\times G_0$-equivariantly upstairs in $\widetilde{M_1}$ (i.e., perform the surgery downstairs and lift). The result is a closed manifold $M_1'$ with $\pi_1(M_1')=\pi_1(M_1)=\Gamma$ (surgery in dim $\geq 5$ on a class in $H_{\geq 2}$ does not change $\pi_1$) and $H_{2k-1}(\widetilde{M_1'};\mathbb{Q})$ has the class $[V]$ killed; all other rational homology classes are unchanged or annihilated by the same handle attachment.

Set $N':=$ universal cover of $M_1'$ truncated to a regular cover with deck group $G_0$ (i.e., $N'$ is the connected $G_0$-Galois cover of $M_1'$ corresponding to the surjection $\pi_1(M_1')=\Gamma\twoheadrightarrow\Gamma/\Gamma_0=G_0$). Then:
\begin{itemize}[leftmargin=2em]
\item $G_0$ acts freely on $N'$ as deck transformations.
\item $\pi_1(N')=\Gamma_0$.
\item $\widetilde{N'}=\widetilde{M_1'}$, which is $\mathbb{Q}$-acyclic by construction.
\item $N'/G_0=M_1'$, a closed smooth manifold.
\end{itemize}

This is the manifold required by the lemma.
\end{proof}

\section{Final Assembly: The Manifold $M$}\label{sec:final}

\begin{theorem}\label{thm:main}
There exists a closed manifold $M$ (without boundary) such that
\begin{enumerate}[label=(\roman*)]
\item $\pi_1(M)\cong\Gamma$, where $\Gamma=\mathrm{SO}(f,\mathcal{O}_F)$ is a uniform lattice in $\mathrm{SO}(5,1)$ (a real semisimple Lie group),
\item $\Gamma$ contains a nontrivial element of order $2$ (the matrix $\gamma$ of Lemma~\ref{lem:torsion}),
\item the universal cover $\widetilde{M}$ is acyclic over $\mathbb{Q}$, i.e., $H_k(\widetilde{M};\mathbb{Q})=0$ for all $k\geq 1$.
\end{enumerate}
\end{theorem}

\begin{proof}
Let $N'$ be the closed smooth $G_0$-manifold from Lemma~\ref{lem:freeaction}, on which $G_0=\Gamma/\Gamma_0$ acts freely. Define
\[
M \;:=\; N'\big/G_0.
\]

\medskip

\noindent\emph{Step 1: $M$ is a closed manifold without boundary.} The action of $G_0$ on $N'$ is smooth and free; for a finite group acting smoothly and freely on a smooth manifold, the quotient is itself a smooth manifold and the projection $N'\to M$ is a smooth covering map. Since $N'$ is compact and without boundary, so is $M$.

\medskip

\noindent\emph{Step 2: Computing $\pi_1(M)$.} The covering $N'\to M$ is regular with deck group $G_0$. The long exact sequence of the fibration
\[
N' \;\to\; M \;\to\; BG_0 \quad\text{(Borel-style classifying space)}
\]
or equivalently the standard exact sequence for a regular covering, gives
\[
1 \;\longrightarrow\; \pi_1(N')\;\longrightarrow\; \pi_1(M) \;\longrightarrow\; G_0 \;\longrightarrow\; 1.
\]
By Lemma~\ref{lem:freeaction}, $\pi_1(N')=\Gamma_0$. Hence we have a short exact sequence
\begin{equation}\label{eq:ses-N}
1 \;\longrightarrow\; \Gamma_0 \;\longrightarrow\; \pi_1(M) \;\longrightarrow\; G_0 \;\longrightarrow\; 1.
\end{equation}

We now identify this extension with the original short exact sequence
\begin{equation}\label{eq:ses-Gamma}
1 \;\longrightarrow\; \Gamma_0 \;\longrightarrow\; \Gamma \;\longrightarrow\; G_0 \;\longrightarrow\; 1.
\end{equation}
Both extensions \eqref{eq:ses-N} and \eqref{eq:ses-Gamma} have the same kernel $\Gamma_0$ and the same quotient $G_0$. The extension class of an extension of $G_0$ by $\Gamma_0$ is determined by the induced outer action $G_0\to\mathrm{Out}(\Gamma_0)$ together with an element of $H^2(G_0;Z(\Gamma_0))$ (Schreier theory), where $Z(\Gamma_0)$ is the center of $\Gamma_0$.

\textbf{Outer action match:} The outer action of $G_0$ on $\Gamma_0$ in \eqref{eq:ses-Gamma} is by definition: pick lift $g'\in\Gamma$ of $g\in G_0$; conjugate by $g'$; this is well-defined modulo $\mathrm{Inn}(\Gamma_0)$. The outer action of $G_0$ on $\Gamma_0$ in \eqref{eq:ses-N} is induced from the deck transformation action of $G_0$ on $N'$, which (by Theorem~\ref{thm:Avramidi}(c) and Theorem~\ref{thm:BQ}) lifts the original action of $G_0$ on $\Gamma_0$ via outer automorphisms. Hence the two outer actions coincide.

\textbf{Center vanishes:} $\Gamma_0$ is a torsion-free uniform lattice in $\mathrm{SO}^\circ(5,1)$, which is a connected simple non-compact Lie group with trivial center (the center of $\mathrm{SO}^\circ(n,1)$ is trivial for $n\geq 2$, see, e.g., Helgason, \emph{Differential Geometry, Lie Groups, and Symmetric Spaces}, Ch.~X). By Borel's density theorem (see Raghunathan, \emph{Discrete Subgroups of Lie Groups}, Theorem 5.5), the centralizer in $\mathrm{SO}^\circ(5,1)$ of $\Gamma_0$ is the center of $\mathrm{SO}^\circ(5,1)$, which is trivial. In particular, $Z(\Gamma_0)\subset C_{\mathrm{SO}^\circ(5,1)}(\Gamma_0)=\{1\}$, so $Z(\Gamma_0)=\{1\}$.

Since $Z(\Gamma_0)=0$, $H^2(G_0;Z(\Gamma_0))=0$, so the extension class in \eqref{eq:ses-N} is determined entirely by the outer action; and similarly for \eqref{eq:ses-Gamma}. The two outer actions agree, so the two extensions are isomorphic, and $\pi_1(M)\cong\Gamma$.

\medskip

\noindent\emph{Step 3: $\widetilde{M}$ is $\mathbb{Q}$-acyclic.} We show $\widetilde{M}=\widetilde{N'}$. The covering $N'\to M$ is finite (with deck group $G_0$). The composition $\widetilde{N'}\to N'\to M$ is a covering map (composition of coverings is a covering when the base has a universal cover, equivalently when $M$ is locally simply connected, which holds for any manifold). And $\widetilde{N'}$ is simply connected. So $\widetilde{N'}$ is the universal cover of $M$:
\[
\widetilde{M} \;=\; \widetilde{N'}.
\]
By Lemma~\ref{lem:freeaction}, $H_k(\widetilde{N'};\mathbb{Q})=0$ for $k\geq 1$, hence $H_k(\widetilde{M};\mathbb{Q})=0$ for $k\geq 1$.

\medskip

\noindent\emph{Step 4: $\Gamma$ has $2$-torsion.} By Lemma~\ref{lem:torsion}, the element $\gamma=\mathrm{diag}(-1,-1,1,1,1,1)\in\Gamma$ has order exactly $2$.

\medskip

\noindent\emph{Step 5: $\Gamma$ is a uniform lattice in a real semisimple Lie group.} By Lemma~\ref{lem:cocompact}, $\Gamma$ is a uniform lattice in $\mathrm{SO}(5,1)$. The Lie group $\mathrm{SO}(5,1)$ is real semisimple: it has Lie algebra $\mathfrak{so}(5,1)$, which is a non-compact real form of the complex simple Lie algebra $\mathfrak{so}(6,\mathbb{C})\cong\mathfrak{sl}(4,\mathbb{C})$ of type $A_3$, hence simple, in particular semisimple.

\medskip

This completes the proof: $M$ is a closed manifold without boundary with $\pi_1(M)\cong\Gamma$ a uniform lattice in the real semisimple Lie group $\mathrm{SO}(5,1)$ containing $2$-torsion, and $\widetilde{M}$ is $\mathbb{Q}$-acyclic.
\end{proof}

\section{Verification of Boundary and Degenerate Cases}\label{sec:edge}

We verify that the construction handles all relevant edge cases, completeness-checking the proof against the gap-audit catalogue.

\medskip

\noindent\textbf{(E1) Existence of Avramidi-style $N$ in dimension $5$.} Theorem~\ref{thm:Avramidi} requires $n\geq 2$ for the base; we have $n=5\geq 2$, so the hypothesis is satisfied. The fiber dimension $d$ may be taken arbitrarily large by tensoring representations or stabilizing the fiber, so we may assume $\dim N = 5+d \geq 8$, satisfying the surgery dimension hypothesis $\dim N\geq 5$ of Theorem~\ref{thm:BQ}.

\medskip

\noindent\textbf{(E2) Free $G_0$-action via stabilization.} The freeness of the $G_0$-action on $N_1=N\times V$ in Lemma~\ref{lem:freeaction} is achieved unconditionally by stabilization with a free $G_0$-sphere $V$, avoiding any need to verify a codimension hypothesis on the original $N$. Hence Theorem~\ref{thm:BQ} is invoked only in its consequence ``free action $\Rightarrow$ ordinary surgery downstairs,'' i.e., we use only Wall's surgery on the quotient $M_1=N_1/G_0$.

\medskip

\noindent\textbf{(E3) Free $G_0$-action exists.} The group $G_0$ is finite (as $[\Gamma:\Gamma_0]<\infty$). Any finite group acts freely on a sphere of sufficiently high dimension: embed $G_0$ in some unitary group $U(n)$ via a faithful unitary representation, which acts freely on $S^{2n-1}\subset\mathbb{C}^n$ outside the origin; restrict the embedding to act freely on $S^{2n-1}$ (this requires the representation to have no nonzero fixed vectors for any nontrivial element, which holds for the regular representation of $G_0$). Thus the auxiliary manifold $V=S^{2n-1}$ in the proof of Lemma~\ref{lem:freeaction} exists.

\medskip

\noindent\textbf{(E4) Connected components of $\mathrm{SO}(5,1)$.} The Lie group $\mathrm{SO}(5,1)$ has two connected components; its identity component $\mathrm{SO}^\circ(5,1)$ has index $2$. Replacing $\Gamma$ by $\Gamma\cap\mathrm{SO}^\circ(5,1)$ if necessary changes $\Gamma$ by at most an index-$2$ subgroup. The element $\gamma=\mathrm{diag}(-1,-1,1,1,1,1)$ has $\det=+1$ and (after computing its action on the standard time-orientation of the form $f$ at $\sigma_1$) preserves time-orientation: at the embedding $\sigma_1$, the negative-definite direction corresponds to the last coordinate $x_6$, on which $\gamma$ acts as $+1$. Hence $\gamma\in\mathrm{SO}^\circ(5,1)$ as well, so $\gamma$ remains a $2$-torsion element after restriction to the identity component if needed.

\medskip

\noindent\textbf{(E5) The case $G_0$ trivial.} If $G_0=\Gamma/\Gamma_0$ were trivial, the image of $\gamma$ in $G_0$ would be trivial, contradicting Lemma~\ref{lem:Gamma0}; so $G_0$ is nontrivial. (We don't need this case, but its absence confirms our construction is non-degenerate.)

\medskip

\noindent\textbf{(E6) The case $\Gamma_0$ trivial.} Impossible: $\Gamma_0$ has finite index in $\Gamma$, but $\Gamma$ is infinite (it is a uniform lattice in a non-compact Lie group, hence infinite).

\medskip

\noindent\textbf{(E7) Center of $\Gamma_0$.} The argument in Step~2 uses $Z(\Gamma_0)=1$. This holds: $\Gamma_0$ is a torsion-free uniform lattice in $\mathrm{SO}^\circ(5,1)$, a center-free Lie group; by Borel density, the Zariski closure of $\Gamma_0$ is the entire group $\mathrm{SO}^\circ(5,1)$ (after complexification), and the centralizer of $\Gamma_0$ in $\mathrm{SO}^\circ(5,1)$ is therefore the center, which is trivial.

\medskip

\noindent\textbf{(E8) Real semisimple Lie group.} The problem asks about a uniform lattice in a real semisimple Lie group. $\mathrm{SO}(5,1)$ is real and semisimple: its Lie algebra $\mathfrak{so}(5,1)$ is a real form of the complex simple Lie algebra of type $D_3$, equivalently $A_3$. So $\Gamma\subset\mathrm{SO}(5,1)$ is indeed a uniform lattice in a real semisimple Lie group.

\section{Discussion}

\begin{remark}[Why hyperbolic?]
The use of a hyperbolic lattice $\Gamma\subset\mathrm{SO}(5,1)$ ensures: (i) cocompactness via Godement's criterion (anisotropic over $F$); (ii) torsion at small primes ($2$-torsion via the diagonal sign-flip element); (iii) PD$_n$ structure of the torsion-free finite-index subgroup. The dimension $5$ is chosen for convenience: any odd dimension $\geq 3$ works, and $5$ gives enough room for the fiber-bundle Avramidi construction to produce a manifold of dimension $\geq 5$ where Browder--Quinn surgery applies.
\end{remark}

\begin{remark}[Why $F=\mathbb{Q}(\sqrt{2})$ and not $\mathbb{Q}$?]
Over $\mathbb{Q}$, an indefinite quadratic form in $\geq 5$ variables (Meyer's theorem) is automatically isotropic, so Godement's criterion fails and the lattice $\mathrm{SO}(f,\mathbb{Z})$ is non-cocompact. The standard remedy in the construction of compact hyperbolic manifolds is to take a totally real number field $F\neq\mathbb{Q}$ and a quadratic form whose Galois conjugates are positive definite, ensuring anisotropy. The simplest such field is $\mathbb{Q}(\sqrt{2})$.
\end{remark}

\begin{remark}[Why $\widetilde{M}$ $\mathbb{Q}$-acyclic but not contractible?]
If $\widetilde{M}$ were contractible, then $M$ itself would be aspherical, and Poincar\'e duality would force $\Gamma$ to be a PD$_{\dim M}$ group; but $\Gamma$ has torsion, and a PD-group is necessarily torsion-free. So one cannot have $\widetilde{M}$ contractible. The relaxation to $\mathbb{Q}$-acyclicity is essential and is exactly what Avramidi's construction delivers.
\end{remark}

\begin{remark}[Connection to the Whitehead--Hattori--Stallings problem]
The construction here is in the spirit of the question whether finite-index extensions of PD$_n$ groups by finite groups can be realized as fundamental groups of closed manifolds. Avramidi's theorem provides the manifold realization, and equivariant surgery promotes a non-free action to a free one.
\end{remark}

\section*{Summary}

\textbf{Answer: YES.} The uniform lattice $\Gamma=\mathrm{SO}(f,\mathcal{O}_F)\subset\mathrm{SO}(5,1)$ for $f=x_1^2+\dots+x_5^2-\sqrt{2}\,x_6^2$ over $F=\mathbb{Q}(\sqrt{2})$ is a uniform lattice in the real semisimple Lie group $\mathrm{SO}(5,1)$, contains the order-$2$ element $\gamma=\mathrm{diag}(-1,-1,1,1,1,1)$, and is the fundamental group of the closed manifold $M=N'/G_0$ constructed via Avramidi's theorem and Browder--Quinn equivariant surgery, whose universal cover is $\mathbb{Q}$-acyclic.

\end{document}
